% GENERATED FILE. Edit canonical lessons, then run npm run generate:books.
% canonical-source-hash: fnv1a64:105d23ea0d168316
% canonical-lessons: HI-C484-datun, HI-C484-katvana, HI-C484-nahlana, HI-C484-malna, HI-C484-ustara

\chapter{Getting Ready}
\label{ch:hi-getting-ready}
\hlchaptermodality{484}

\begin{chapteropening}
\textbf{By the end of this chapter:} \emph{I can talk about washing and grooming.}

Complete the last lesson of chapter 484: I can talk about washing and grooming.
\end{chapteropening}

\section[\texorpdfstring{\hi{दातुन} (dātun)}{दातुन (dātun)}]{\hi{दातुन} (dātun) — a twig toothbrush}

\label{lesson:HI-C484-datun}



\begin{quote}
Before the new one: say the Hindi for force, then the Hindi for to tire someone out.
\end{quote}

\subsection*{You'll want to know: \hi{दातुन}}
\textbf{\hi{दातुन}} — \emph{dātun} — \textquotedblleft{}a twig toothbrush\textquotedblright{}.

\begin{quote}
\hi{गाँव} \hi{में} \hi{दादा} \hi{आज} \hi{भी} \hi{दातुन} \hi{से} \hi{दाँत} \hi{साफ़} \hi{करते} \hi{हैं।}
\end{quote}

\begin{quote}
\emph{gā̃v meṁ dādā āj bhī dātun se dā̃t sāf karte haiṁ.}
\end{quote}

\begin{quote}
In the village my grandfather still cleans his teeth with a twig.
\end{quote}

Feminine.

\begin{cousinweb}
\textbf{\hi{दातुन}} (\emph{dātun}) is a tadbhava: CDIAL lists Hindi \emph{dātun}, \textquotedblleft{}tooth stick\textquotedblright{}, under Sanskrit \emph{dantapavana}, \textquotedblleft{}tooth-cleaning\textquotedblright{} (CDIAL 6157): a twig, often of neem, chewed at one end.
\end{cousinweb}

\begin{grammarlens}[title={What you've built}]
One more word for describing what you lived and felt.
\end{grammarlens}

\subsection*{Your turn}
\begin{itemize}
\raggedright
  \item \emph{Say it:} \emph{dātun}
  \item \emph{Say it:} \emph{dātun}, once more
  \item \emph{Say it:} \emph{gā̃v meṁ dādā āj bhī dātun se dā̃t sāf karte haiṁ.} — In the village my grandfather still cleans his teeth with a twig.
  \item \emph{Recall:} say the Hindi for force, then the Hindi for to tire someone out, then say \emph{dātun} again
\end{itemize}

\begin{tcolorbox}[breakable,colback=teal!4,colframe=teal!35!black,title={Before you move on}]
What does \hi{दातुन} mean? (\textquotedblleft{}A twig toothbrush\textquotedblright{}.) Say it once more.
\end{tcolorbox}
\section[\texorpdfstring{\hi{कटवाना} (kaṭvānā)}{कटवाना (kaṭvānā)}]{\hi{कटवाना} (kaṭvānā) — to have (something) cut}

\label{lesson:HI-C484-katvana}



\begin{quote}
Before the new one: say the Hindi for to tire someone out, then the Hindi for a twig toothbrush.
\end{quote}

\subsection*{You'll want to know: \hi{कटवाना}}
\textbf{\hi{कटवाना}} — \emph{kaṭvānā} — \textquotedblleft{}to have (something) cut\textquotedblright{}.

\begin{quote}
\hi{शादी} \hi{से} \hi{पहले} \hi{अरुण} \hi{ने} \hi{नाई} \hi{से} \hi{बाल} \hi{कटवाए।}
\end{quote}

\begin{quote}
\emph{śādī se pahle Aruṇ ne nāī se bāl kaṭvāe.}
\end{quote}

\begin{quote}
Before the wedding Arun had his hair cut at the barber's.
\end{quote}

Past, with \hi{ने} (he or she …): \hi{उसने} \hi{बाल} \hi{कटवाए}, \emph{usne bāl kaṭvāe}.

\begin{cousinweb}
\textbf{\hi{कटवाना}} (\emph{kaṭvānā}) is the causative of \emph{kāṭnā}, \textquotedblleft{}to cut\textquotedblright{}, which CDIAL lists under Sanskrit \emph{kártati}, \textquotedblleft{}cuts\textquotedblright{} (CDIAL 2854): with \emph{-vā-}, someone else does the cutting for you. The person who does it takes \emph{se}: \emph{nāī se}.
\end{cousinweb}

\begin{grammarlens}[title={What you've built}]
One more word for describing what you lived and felt.
\end{grammarlens}

\subsection*{Your turn}
\begin{itemize}
\raggedright
  \item \emph{Say it:} \emph{kaṭvānā}
  \item \emph{Say it:} \emph{kaṭvānā}, once more
  \item \emph{Say it:} \emph{śādī se pahle Aruṇ ne nāī se bāl kaṭvāe.} — Before the wedding Arun had his hair cut at the barber's.
  \item \emph{Recall:} say the Hindi for to tire someone out, then the Hindi for a twig toothbrush, then say \emph{kaṭvānā} again
\end{itemize}

\begin{tcolorbox}[breakable,colback=teal!4,colframe=teal!35!black,title={Before you move on}]
What does \hi{कटवाना} mean? (\textquotedblleft{}To have (something) cut\textquotedblright{}.) Say it once more.
\end{tcolorbox}
\section[\texorpdfstring{\hi{नहलाना} (nahlānā)}{नहलाना (nahlānā)}]{\hi{नहलाना} (nahlānā) — to bathe someone}

\label{lesson:HI-C484-nahlana}



\begin{quote}
Before the new one: say the Hindi for a twig toothbrush, then the Hindi for to have (something) cut.
\end{quote}

\subsection*{You'll want to know: \hi{नहलाना}}
\textbf{\hi{नहलाना}} — \emph{nahlānā} — \textquotedblleft{}to bathe someone\textquotedblright{}.

\begin{quote}
\hi{माँ} \hi{रोज़} \hi{सुबह} \hi{बच्चे} \hi{को} \hi{नहलाती} \hi{है।}
\end{quote}

\begin{quote}
\emph{māṁ roz subah bacce ko nahlātī hai.}
\end{quote}

\begin{quote}
Mother bathes the baby every morning.
\end{quote}

Past, with \hi{ने} (he or she …): \hi{उसने} \hi{बच्चे} \hi{को} \hi{नहलाया}, \emph{usne bacce ko nahlāyā}.

\begin{cousinweb}
\textbf{\hi{नहलाना}} (\emph{nahlānā}) is the causative partner of \emph{nahānā}, \textquotedblleft{}to bathe\textquotedblright{}, which CDIAL lists under Sanskrit \emph{snāti}, \textquotedblleft{}bathes\textquotedblright{} (CDIAL 13786): to give someone else a bath.
\end{cousinweb}

\begin{grammarlens}[title={What you've built}]
One more word for describing what you lived and felt.
\end{grammarlens}

\subsection*{Your turn}
\begin{itemize}
\raggedright
  \item \emph{Say it:} \emph{nahlānā}
  \item \emph{Say it:} \emph{nahlānā}, once more
  \item \emph{Say it:} \emph{māṁ roz subah bacce ko nahlātī hai.} — Mother bathes the baby every morning.
  \item \emph{Recall:} say the Hindi for a twig toothbrush, then the Hindi for to have (something) cut, then say \emph{nahlānā} again
\end{itemize}

\begin{tcolorbox}[breakable,colback=teal!4,colframe=teal!35!black,title={Before you move on}]
What does \hi{नहलाना} mean? (\textquotedblleft{}To bathe someone\textquotedblright{}.) Say it once more.
\end{tcolorbox}
\section[\texorpdfstring{\hi{मलना} (malnā)}{मलना (malnā)}]{\hi{मलना} (malnā) — to rub in; to massage}

\label{lesson:HI-C484-malna}



\begin{quote}
Before the new one: say the Hindi for to have (something) cut, then the Hindi for to bathe someone.
\end{quote}

\subsection*{You'll want to know: \hi{मलना}}
\textbf{\hi{मलना}} — \emph{malnā} — \textquotedblleft{}to rub in; to massage\textquotedblright{}.

\begin{quote}
\hi{नहाने} \hi{से} \hi{पहले} \hi{दादी} \hi{ने} \hi{बच्चे} \hi{के} \hi{पैरों} \hi{पर} \hi{तेल} \hi{मला।}
\end{quote}

\begin{quote}
\emph{nahāne se pahle dādī ne bacce ke pairoṁ par tel malā.}
\end{quote}

\begin{quote}
Before the bath, my grandmother rubbed oil into the baby's legs.
\end{quote}

Past, with \hi{ने} (he or she …): \hi{उसने} \hi{हाथों} \hi{पर} \hi{तेल} \hi{मला}, \emph{usne hāthoṁ par tel malā}.

\begin{cousinweb}
\textbf{\hi{मलना}} (\emph{malnā}) is a tadbhava: CDIAL lists Hindi \emph{malnā}, \textquotedblleft{}to rub, shampoo\textquotedblright{}, under \emph{*marati}, a form it reconstructs (CDIAL 9870). An oil massage for a baby is \emph{māliś}, and the verb for it is \emph{malnā}.
\end{cousinweb}

\begin{grammarlens}[title={What you've built}]
One more word for describing what you lived and felt.
\end{grammarlens}

\subsection*{Your turn}
\begin{itemize}
\raggedright
  \item \emph{Say it:} \emph{malnā}
  \item \emph{Say it:} \emph{malnā}, once more
  \item \emph{Say it:} \emph{nahāne se pahle dādī ne bacce ke pairoṁ par tel malā.} — Before the bath, my grandmother rubbed oil into the baby's legs.
  \item \emph{Recall:} say the Hindi for to have (something) cut, then the Hindi for to bathe someone, then say \emph{malnā} again
\end{itemize}

\begin{tcolorbox}[breakable,colback=teal!4,colframe=teal!35!black,title={Before you move on}]
What does \hi{मलना} mean? (\textquotedblleft{}To rub in; to massage\textquotedblright{}.) Say it once more.
\end{tcolorbox}
\section[\texorpdfstring{\hi{उस्तरा} (ustarā)}{उस्तरा (ustarā)}]{\hi{उस्तरा} (ustarā) — a razor}

\label{lesson:HI-C484-ustara}



\begin{quote}
Before the new one: say the Hindi for to bathe someone, then the Hindi for to rub in.
\end{quote}

\subsection*{You'll want to know: \hi{उस्तरा}}
\textbf{\hi{उस्तरा}} — \emph{ustarā} — \textquotedblleft{}a razor\textquotedblright{}.

\begin{quote}
\hi{नाई} \hi{ने} \hi{उस्तरे} \hi{से} \hi{दाढ़ी} \hi{बनाई।}
\end{quote}

\begin{quote}
\emph{nāī ne ustare se dāṛhī banāī.}
\end{quote}

\begin{quote}
The barber shaved the beard with a razor.
\end{quote}

Masculine.

\begin{cousinweb}
\textbf{\hi{उस्तरा}} (\emph{ustarā}) is usually traced to Persian \emph{ustura}, \textquotedblleft{}razor\textquotedblright{}. \emph{dāṛhī banānā}, \textquotedblleft{}to make the beard\textquotedblright{}, is how Hindi says to shave.
\end{cousinweb}

\begin{grammarlens}[title={What you've built}]
One more word for describing what you lived and felt.
\end{grammarlens}

\subsection*{Your turn}
\begin{itemize}
\raggedright
  \item \emph{Say it:} \emph{ustarā}
  \item \emph{Say it:} \emph{ustarā}, once more
  \item \emph{Say it:} \emph{nāī ne ustare se dāṛhī banāī.} — The barber shaved the beard with a razor.
  \item \emph{Recall:} say the Hindi for to bathe someone, then the Hindi for to rub in, then say \emph{ustarā} again
\end{itemize}

\begin{tcolorbox}[breakable,colback=teal!4,colframe=teal!35!black,title={Before you move on}]
What does \hi{उस्तरा} mean? (\textquotedblleft{}A razor\textquotedblright{}.) Say it once more.
\end{tcolorbox}
